% VERITAS: Independence-Weighted Multi-Model Consensus
% for AI Output Verification
%
% Preprint — submitted to arXiv cs.AI March 2026
% Source: https://github.com/RJLopezAI/veritas
%
% This file: papers/veritas_mis_greedy_2026.tex

\documentclass[11pt]{article}
\usepackage{hyperref}
\usepackage{booktabs}
\usepackage{amsmath}

\title{VERITAS: Independence-Weighted Multi-Model Consensus\\
       for AI Output Verification}

\author{
  RJ Lopez\\
  AegisAudits\\
  \texttt{api@aegisaudits.com}
}

\date{March 2026}

\begin{document}
\maketitle

\begin{abstract}
We present VERITAS, a production architecture for hallucination detection
in AI agent pipelines. The core contribution is MIS\_GREEDY
(Maximally Independent Set, Greedy), an independence-weighted voting
algorithm that addresses the structural limitation of single-model
hallucination detectors: a model cannot detect its own mode collapse.

VERITAS submits claims to three language models simultaneously
(gemini-2.0-flash-001, gemini-1.5-pro-002, gemini-2.0-flash-thinking),
applies MIS\_GREEDY independence scoring, and returns a structured
verdict with confidence score, consensus flag, and concern flags.

We evaluate against Galileo's Luna-2 \cite{galileo2026luna2},
the current state-of-the-art single-model SLM approach reporting
88\% accuracy at 152ms average latency. On a 35-claim benchmark
targeting confidence-trap claims --- cases where single-model training
distributions contain dense, confident, incorrect associations ---
VERITAS achieves 97.1\% overall accuracy and 100\% on legal, medical,
and financial domain claims at 59ms average latency.

The system is deployed on Google Cloud Run with Vertex AI inference
and BigQuery audit logging. Source code, benchmark dataset, and
OpenAPI specification are available at
\url{https://github.com/RJLopezAI/veritas}.
\end{abstract}

\section{Introduction}

AI agent pipelines that return factual claims to users require a
verification gate before output. The naive approach --- asking the
generating model to verify its own output --- fails on the class of
claims most likely to cause harm: those where the model's training
distribution contains a confident, dense, incorrect association.

We term these \textit{confidence-trap claims}. Examples include legal
myths that appear thousands of times as fact in training corpora
(``a verbal contract is unenforceable''), medical misconceptions
reinforced by pop science content (``sugar causes hyperactivity''),
and financial misunderstandings treated as settled knowledge
(``an LLC provides absolute personal liability protection'').

A single model fine-tuned for hallucination detection --- the approach
taken by Galileo's Luna-2 \cite{galileo2026luna2} --- improves accuracy
on the general distribution but cannot escape confident wrong answers
baked into its base model's weights.

Multi-model consensus addresses this structurally. Three models with
different architectures and training histories are unlikely to share
the same mode collapse. When they disagree, that disagreement is
the detection signal.

\section{The MIS\_GREEDY Algorithm}

Standard ensemble voting treats unanimous agreement as maximally strong
evidence. MIS\_GREEDY inverts this for the specific case of mode collapse
detection: suspicious unanimous agreement (low confidence spread) is
treated as weaker evidence than a healthy 2-1 split.

\begin{equation}
w = \begin{cases}
0.75 & \text{if } |\text{verdicts}| = 1 \text{ and spread} < 0.05 \\
0.90 & \text{if } |\text{verdicts}| = 1 \text{ and spread} \geq 0.05 \\
1.00 & \text{if } |\text{verdicts}| = 2 \text{ (healthy disagreement)} \\
0.85 & \text{if } |\text{verdicts}| = 3 \text{ (three-way split)}
\end{cases}
\end{equation}

where spread $= \max(\text{confidences}) - \min(\text{confidences})$.

The final confidence score is the weighted agreement score:
$\hat{c} = \text{agreement}(votes) \times w$

Consensus is declared when $\hat{c} \geq 0.80$ (standard) or
$\hat{c} \geq 0.90$ (legal, medical domains).

\section{Architecture}

Three Vertex AI model calls are issued in parallel via
\texttt{asyncio.gather}. Wall-clock latency equals the slowest
model call rather than their sum, producing the observed 59ms
average latency despite three concurrent inference calls.

Every verification call is logged to BigQuery with: claim hash,
per-model verdicts, confidence scores, MIS\_GREEDY independence weight,
final consensus verdict, domain, and latency. This audit trail
accumulates as a labeled verification dataset growing with production usage.

The system is exposed as both a REST API (OpenAPI 3.1 specification
at \url{https://veritas-toll-road-m72j3qteca-uc.a.run.app/openapi.yaml}) and an MCP
tool server compatible with Claude Desktop and any MCP-capable
agent runtime.

\section{Evaluation}

\subsection{Benchmark Design}

We construct a 35-claim benchmark across three categories:

\begin{itemize}
\item \textbf{Category A} (plausible falsehoods, $n=8$): Claims appearing
densely in training corpora in misleading form.
\item \textbf{Category B} (high-stakes domain claims, $n=22$): Legal,
medical, and financial claims with confidence-trap properties.
\item \textbf{Category C} (recency collision, $n=5$): Claims accurate
in training data but since changed.
\end{itemize}

Ground truth labels are verified against primary sources.
The benchmark script and full claim set are publicly available
at \url{https://github.com/RJLopezAI/veritas/tree/main/benchmark/}.

\subsection{Results}

\begin{table}[h]
\centering
\begin{tabular}{lcc}
\toprule
Metric & VERITAS & Luna-2 \\
\midrule
Overall accuracy & \textbf{97.1\%} & 88.0\% \\
Category B (domain) & \textbf{100\%} & $\sim$82\% est. \\
Confidence-trap claims & \textbf{100\%} & $\sim$75\% est. \\
Average latency & \textbf{59ms} & 152ms \\
p95 latency & 162ms & --- \\
Cost per call & \$0.05 & \$0.0002 \\
\bottomrule
\end{tabular}
\caption{VERITAS vs Luna-2. Luna-2 figures from \cite{galileo2026luna2}.
Domain and confidence-trap figures for Luna-2 are estimates; Galileo
has not published domain-specific breakdowns.}
\end{table}

The single failure (C002) occurred on a recency collision claim,
reflecting the shared training cutoff limitation of all model-based
approaches. Zero failures on all 22 Category B claims.

\subsection{Cost Analysis}

VERITAS costs \$0.05/call versus Luna-2's \$0.0002/call (250x differential).
Break-even: one wrong verdict costs more than $\$0.05 / 0.12 = \$0.42$.
In legal, medical, and financial contexts this threshold is routinely exceeded.

\section{Discussion}

The latency result (59ms average vs Luna-2's 152ms) was unexpected.
The \texttt{asyncio.gather} parallel architecture captures wall-clock
time of the fastest model rather than the sum, and Vertex AI's
gemini-2.0-flash inference path produces lower average latency
than Luna-2's hosted SLM infrastructure.

The practical positioning: Luna-2 is correct for high-volume,
low-stakes general guardrailing. VERITAS is correct for the specific
tail of high-stakes domain verification where mode collapse is the
primary failure mode and wrong confident verdicts cause real harm.

\section{Conclusion}

A model cannot detect its own mode collapse. Multi-model consensus
with independence-weighted voting addresses this structurally.
VERITAS achieves 97.1\% overall accuracy and 100\% on confidence-trap
claims in legal, medical, and financial domains, at 59ms average latency,
outperforming single-model SLM approaches on both dimensions relevant
to high-stakes deployment.

\bibliographystyle{plain}
\begin{thebibliography}{1}
\bibitem{galileo2026luna2}
Galileo AI.
\newblock Luna-2: High Accuracy, Low Cost Evals via Fine-Tuned SLMs.
\newblock arXiv preprint arXiv:2602.18583, 2026.
\newblock \url{https://arxiv.org/pdf/2602.18583}
\end{thebibliography}

\end{document}
